% Generated draft: requires mathematical and source review.
\documentclass[11pt]{article}
\usepackage[margin=25mm]{geometry}
\usepackage{fontspec}
\setmainfont{Latin Modern Roman}
\usepackage{amsmath,amssymb,mathtools,unicode-math}
\setmathfont{Latin Modern Math}
\usepackage{enumitem,needspace,xurl,hyperref}
\setlength{\parindent}{0pt}
\setlength{\parskip}{5pt}
\setlength{\emergencystretch}{3em}
\newcommand{\N}{\mathbb{N}}
\newcommand{\Z}{\mathbb{Z}}
\newcommand{\Q}{\mathbb{Q}}
\newcommand{\R}{\mathbb{R}}
\newcommand{\C}{\mathbb{C}}
\newcommand{\F}{\mathbb{F}}
\newcommand{\A}{\mathbb{A}}
\newcommand{\PP}{\mathbb{P}}
\newcommand{\E}{\mathbb{E}}
\newcommand{\eps}{\varepsilon}
\newcommand{\dd}{\,\mathrm{d}}
\begin{document}
\section*{UnsolvedMath ID: 30001628}
\textbf{Source problem code:} catalogue entry
\subsection{Definitions and mathematical statement}
Let \((X,\tau)\) be a non-Hausdorff topological space. Let \(S\subseteq X\) be closed and satisfy the supplied definition of a \emph{closed crack}: for every \(x,y\in X\),
\[
\bigl(x\ne y\ \text{and}\ \{x,y\}\nsubseteq S\bigr)
\Longrightarrow
\exists U,V\in\tau\quad x\in U,\quad y\in V,\quad U\cap V=\varnothing.
\]
Put \(R=X\setminus S\), equipped initially with its subspace topology. The source calls \(R\) the regular part and \(S\) the singular part; these names impose no additional stated hypotheses. Write \(2^S=\mathcal P(S)\) for the power set of \(S\).

The supplied construction question asks how to choose a family \(L\subseteq\mathcal P(S)\) of added points and a topology \(\widetilde\tau\) on the disjoint union
\[
\widetilde X=L\sqcup R
\]
so as to repair the non-Hausdorff quotient referred to in the source.

\textbf{Unresolved formulation.} The supplied extraction does not specify the quotient construction, an equivalence relation, a map relating \(X\) to \(\widetilde X\), or the properties that constitute a successful repair. In particular, it does not explicitly require \(\widetilde X\) to be Hausdorff or specify whether its topology must induce the original subspace topology on \(R\). Thus the requested choice of \(L\) and \(\widetilde\tau\) is a construction question with an unspecified success criterion, rather than a fully determined existence assertion.
\subsection{Short English statement}
Given a non-Hausdorff space whose pairs of distinct points can be separated by disjoint neighborhoods whenever at least one point lies outside a specified closed singular set \(S\), choose subsets of \(S\) as replacement points and topologize their disjoint union with \(X\setminus S\). The intended meaning of a successful repair is missing from the supplied extraction.
\subsection{Sources}
\begin{itemize}
\item Supplied local UnsolvedMath registry record, permanent ID \(30001628\), title \emph{Topological Repair of Non-Hausdorff Quotient Spaces}, source problem code \texttt{catalogue entry}, in the Oberwolfach Reports Open Problems collection. No individual proposer is identified; the supplied proposed year is \(2010\).
\item Supplied report citation: \emph{Representation Theory and Harmonic Analysis} (2011), Oberwolfach Reports, pp.\ 3043--3083, DOI \texttt{10.4171/owr/2010/52}. Source URL and literature-source link labeled ``original OWR report'': \url{https://doi.org/10.4171/owr/2010/52}. The supplied background describes Oberwolfach Reports as documenting mathematical workshops and their open-problem discussions.
\item Literature assessment dated \texttt{2026-08-21}: supplied dated curation, not independently verified here. It labels the problem open, reports that the defining context of the singular quotient construction is incomplete in the extraction, and states that the checked report and literature do not identify a settled canonical topology for the stated object. The supplied curation requests independent verification of time-sensitive status and cited claims.
\item The registry labels the statement \texttt{corrected\_verified} and says it was checked against the official report, preserving the hypotheses, notation, and open target. This is a supplied verification claim, not an independent verification of the report or of mathematical correctness in this draft.
\end{itemize}
\end{document}
